\chapter{前言}
网络爬虫的核心工作不是“下载一个网页”，而是把来源、请求、解析和数据质量连成一条可靠的链。页面结构会变化，链接可能跳转，字段可能缺失；能运行一次的脚本，还算不上可复用的采集程序。

本书按“问题—原理—代码—验证”讲解。四部分依次处理基础链路、采集进阶、专题边界与工程实践。XPath 紧接 HTML 解析；浏览器、Scrapy 和异步客户端在基础链路之后登场；AST 与分布式放在后段。这样读到新工具时，先知道它要解决什么问题。

每章列出前置知识、输入与输出，以及一项可核对的验收标准。代码保存在独立文件，依赖版本、运行命令和期望输出见附录的复现索引；项目中的自动化测试覆盖成功路径和常见失败路径。示例优先使用固定样本或本地 HTTP 服务，不把教学效果押在外站是否在线上。

采集真实数据前，应确认站点规则、接口约定、账户权限和数据用途，控制请求速率，避免收集不需要的个人信息。\texttt{robots.txt} 不是授权凭证；代理、验证码、证书固定和付费墙也不是默认要“攻克”的关卡。遇到边界不清，先停下来确认，通常比写出更长的脚本划算。
